\documentclass[noauthor,noinstructornotes]{ximera}

\title{Week 5}
\author{Math 1193 Design Team}

\begin{document}
\begin{abstract}
    Activities for Week 5.
\end{abstract}
\maketitle

\begin{problem}
There is a special point on the graph of a parabola called the 
\textbf{vertex}. For the parent function \(f\left(x\right)=x^2\), 
the vertex is located at the point \((0,0)\). Each transformation 
applied to the parent function moves the vertex.

\begin{enumerate}
    \item What is the vertex of the graph of \(g(x) = (x - 2)^2\)?
    \item What is the vertex of the graph of \(g(x) = x^2 - 2\)?
    \item What is the vertex of the graph of \(g(x) = 2x^2\)?
    \item What is the vertex of the graph of \(g(x) = (-x)^2\)?
\end{enumerate}

\begin{instructorNotes}
You might take the opportunity here to show a demo on Desmos to review
how transformations work.
\end{instructorNotes}
\end{problem}

\begin{problem}
\begin{enumerate}
    \item Write down the equation of a parabola whose graph has
    its vertex at \((-1, 3)\).
    \item Can you find another?
    \item How many such parabolas are there?
\end{enumerate}

\begin{instructorNotes}
If students have difficulty on \verb+vertexTransformations+, it may 
be worth it to spend extra time here. 

This could also be skipped.
\end{instructorNotes}
\end{problem}

\begin{problem}
A quadratic function is one of the form \(ax^2 + bx + c\), where \(a\),
\(b\), and \(c\) are real numbers (some may be 0). Select all the 
following functions which are quadratic. 

For each quadratic function below, find its zeros (zeros are 
\(x\)-values where the functions output is 0). Check with your group.
Did you all use the same method? Can you understand each other's methods?

\begin{enumerate}
    \item \(g\left(x\right)=x^2-9\)
    \item \(f\left(x\right)=6+5x+x^2\)
    \item \(g\left(x\right)=0.1x^2\)
    \item \(h\left(x\right)=\left(x-1\right)\left(x-2\right)^2\)
    \item \(f\left(x\right)=2\left(x+2\right)\left(x-9\right)\)
    \item \(f\left(x\right)=\left(3x-1\right)\left(x+1\right)\)
    \item \(g\left(x\right)=-3\left(x-\frac{1}{2}\right)^2+4\)
    \item \(h\left(x\right)=x\left(x+3\right)\)
    \item \(f\left(x\right)=2\left(x-6\right)\)
    \item \(g\left(x\right)=2x^2-11x+5\)
\end{enumerate}

\begin{instructorNotes}
Parts will have to be omitted based on time restraints.
\end{instructorNotes}

\end{problem}

\begin{problem}

\begin{enumerate}
    \item Write down a quadratic whose zeros are 2 and 5.
    \item Can you find another?
    \item How many such quadratics are there?
\end{enumerate}

\begin{instructorNotes}
If students are stuck, emphasize it's okay to put down a guess, then
see if it works and if not, how it can be modified.
\end{instructorNotes}

\end{problem}

\begin{problem}
Alana and Grayson were given the following problem:

\textit{Given a function \(f(x) = x^2 + 7\), find \(f(x - 2)\).}

\begin{itemize}
    \item Alana's work is as follows:
    \begin{align*}
        f(x - 2) & = (x - 2)^2 + 7 \\
        & = x^2 - 2^2 + 7 \\
        & = x^2 - 4 + 7 \\
        & = x^2 + 3
    \end{align*}
    \item Grayson's work is as follows:
    \begin{align*}
        f(x - 2) & = (x - 2)^2 + 7 \\
        & = x^2 - 4x + 4 + 7 \\
        & = x^2 - 4x + 11
    \end{align*}
\end{itemize}

Whose work is correct? Have a discussion with a partner. Try to ask
each other why.
\end{problem}

\end{document}
